% !TEX root = ../main.tex
% Figure: CWT upgrade map from 2D minimal model to higher-order/3D model

\begin{figure}[htbp]
  \centering
  \resizebox{0.98\textwidth}{!}{%
  \begin{tikzpicture}[
    node distance=9mm and 8mm,
    box/.style={rectangle, rounded corners, draw, align=center, minimum width=3.4cm, minimum height=1.0cm, fill=#1!15},
    out/.style={rectangle, draw, align=left, minimum width=3.2cm, minimum height=0.9cm, fill=#1!10},
    arr/.style={-Stealth, thick}
  ]
    \node[box=blue] (base) {2D minimal CWT\\(4-wave + one vertical profile)};
    \node[box=teal, below=of base] (q1) {Question: where does it fail?};

    \node[box=green, below left=14mm and 24mm of q1] (v) {Recover DOF-1\\Vertical modal basis};
    \node[box=orange, below=14mm of q1] (r) {Recover DOF-2\\Radiation channels};
    \node[box=purple, below right=14mm and 24mm of q1] (e) {Recover DOF-3\\Finite-size envelope};

    \node[out=green, below=8mm of v] (ov) {$\rightarrow$ Better mode ordering\\$\rightarrow$ Better overlap factors};
    \node[out=orange, below=8mm of r] (or) {$\rightarrow$ Up/down leakage split\\$\rightarrow$ Better polarization/far field};
    \node[out=purple, below=8mm of e] (oe) {$\rightarrow$ Edge leakage captured\\$\rightarrow$ Better array-size prediction};

    \node[box=red, below=16mm of r, minimum width=5.7cm] (target) {Higher-order / 3D CWT output\\usable interfaces for RCWA/FEM/FDTD and device models};

    \draw[arr] (base) -- (q1);
    \draw[arr] (q1) -- (v);
    \draw[arr] (q1) -- (r);
    \draw[arr] (q1) -- (e);
    \draw[arr] (v) -- (ov);
    \draw[arr] (r) -- (or);
    \draw[arr] (e) -- (oe);
    \draw[arr] (ov.south) |- (target.west);
    \draw[arr] (or.south) -- (target.north);
    \draw[arr] (oe.south) |- (target.east);
  \end{tikzpicture}
  }
  \caption{CWT upgrade map used in this chapter: from minimal 2D CWT to higher-order/3D CWT by explicitly recovering three missing degrees of freedom.}
  \label{fig:ch10_cwt_upgrade_map}
\end{figure}
